% ==========================================
% ФАЙЛ: p7.tex
% БЛОК БИЛЕТОВ: Числовые ряды (Признаки сходимости)
% ==========================================

% --- Вопрос 31 ---
\examquestion{31}{Линейные свойства рядов и влияние конечных изменений}
{Сформулируйте и докажите свойства линейности числовых рядов. Докажите также теорему о том, что отбрасывание, изменение или добавление конечного числа членов ряда влияет на значение его суммы, но оставляет абсолютно инвариантным сам факт его сходимости или расходимости.}

\paragraph{База: Что такое ряд и его сумма}
Прежде чем что-то доказывать, зафиксируем базу. Числовой ряд $\sum_{n=1}^\infty a_n$ — это бесконечная сумма. 
\textbf{Частичная сумма} $S_n$ — это сумма первых $n$ слагаемых: $S_n = \sum_{k=1}^n a_k$.
Ряд называется \textbf{сходящимся}, если существует конечный предел его частичных сумм при устремлении $n$ в бесконечность: $\lim_{n \to \infty} S_n = S$. Это число $S$ и есть сумма ряда. Если предела нет или он равен бесконечности — ряд расходится.

\paragraph{Суть на пальцах}
Умножение ряда на число или сложение двух рядов работает точно так же, как с обычными пределами — всё линейно. Что касается изменения конечного числа элементов: представьте бесконечный поезд. Если вы отцепите пару вагонов в начале или поменяете в них груз, общий вес поезда изменится. Но если поезд был бесконечно тяжелым (ряд расходился), он таким и останется. Если его вес был конечным (сходился), то и останется конечным, просто итоговая цифра на весах станет другой.

\paragraph{Свойства линейности}
Если ряды $\sum_{n=1}^\infty a_n$ и $\sum_{n=1}^\infty b_n$ сходятся к конечным суммам $A$ и $B$ соответственно, то:
1. Ряд $\sum_{n=1}^\infty (c \cdot a_n)$ сходится, и его сумма равна $c \cdot A$ (где $c \in \mathbb{R}$).
2. Ряд $\sum_{n=1}^\infty (a_n \pm b_n)$ сходится, и его сумма равна $A \pm B$.

\begin{proof}[Доказательство линейности]
1. По определению сходимости, пределы частичных сумм исходных рядов существуют и конечны: $\lim_{n \to \infty} S_n^A = A$ и $\lim_{n \to \infty} S_n^B = B$.
2. Рассмотрим частичную сумму для новой линейной комбинации рядов: $\sum_{k=1}^n (c \cdot a_k + d \cdot b_k)$.
3. Так как $n$ — это конечное число, мы можем перегруппировать слагаемые и вынести константы по обычным правилам арифметики:
$$ S_n = c \sum_{k=1}^n a_k + d \sum_{k=1}^n b_k = c S_n^A + d S_n^B $$
4. Теперь перейдем к пределу при $n \to \infty$. Воспользуемся классической теоремой матанализа о линейности предела последовательности:
$$ \lim_{n \to \infty} S_n = c \lim_{n \to \infty} S_n^A + d \lim_{n \to \infty} S_n^B = cA + dB $$
Поскольку предел существует и равен конечному числу $cA + dB$, новый ряд сходится.
\end{proof}

\paragraph{Теорема о конечных изменениях}
Отбрасывание, добавление или изменение конечного числа членов ряда не влияет на его сходимость или расходимость (но может изменить итоговую сумму).

\begin{proof}
1. Пусть дан ряд $\sum_{n=1}^\infty a_n$. Допустим, мы изменили или отбросили первые $N$ членов.
2. Рассмотрим частичную сумму исходного ряда $S_n$ для достаточно большого номера $n$ (где $n > N$). Мы можем разбить её на две части — измененную "голову" и нетронутый "хвост":
$$ S_n = \sum_{k=1}^n a_k = \left( \sum_{k=1}^N a_k \right) + \sum_{k=N+1}^n a_k $$
3. Обозначим сумму первых $N$ членов как $C$. Это просто фиксированное конечное число (ведь слагаемых конечное количество). Вторую сумму (частичную сумму хвоста) обозначим как $T_n$. 
Получаем: $S_n = C + T_n$.
4. Выразим хвост: $T_n = S_n - C$.
5. Перейдем к пределу при $n \to \infty$. Константа $C$ не влияет на сам факт существования предела. Предел $\lim_{n \to \infty} S_n$ существует (весь ряд сходится) тогда и только тогда, когда существует предел $\lim_{n \to \infty} T_n$ (хвост сходится).
Следовательно, любые махинации с конечным куском $C$ меняют только итоговое значение, но никак не влияют на сам факт сходимости.
\end{proof}


% --- Для Вопроса 31 (Линейность и конечные изменения) ---
\begin{tcolorbox}[colback=gray!5, colframe=black, boxrule=0.5pt, arc=0pt]
\textbf{Резюмируем:} \\
\textbf{Тезис:} Линейные комбинации сходятся ($cA+dB$). Изменение конечного числа членов не влияет на сходимость. \\
\textbf{Ход (Линейность):} $S_n = c S_n^A + d S_n^B$. Берем предел, по т. о линейности предела получаем $cA+dB$. \\
\textbf{Ход (Хвост):} $S_n = C + T_n$, где $C$ это конечная сумма. При $n \to \infty$ константа $C$ не влияет на $\lim S_n$. Сходимость зависит только от хвоста $T_n$.
\end{tcolorbox}




% --- Вопрос 32 ---
\examquestion{32}{Ряды с неотрицательными членами и признаки сравнения}
{Сформулируйте и докажите критерий сходимости неотрицательных рядов через ограниченность сверху частичных сумм. Проведите доказательство признаков сравнения в форме неравенств и в предельной форме. Покажите на конкретном примере, как предельный признак сравнения позволяет исследовать на сходимость ряд, общий член которого содержит тригонометрическую или логарифмическую зависимость.}

\paragraph{Суть на пальцах}
Если вы складываете только положительные числа (или нули), ваша итоговая сумма может только расти. У неё есть два пути: либо она упрется в какой-то невидимый потолок (сходится к числу), либо улетит в бесконечность (расходится). Колебаться она не может. Признак сравнения — это логика: если «большой» ряд сходится, то и «меньший» подавно. А предельный признак позволяет откинуть сложный мусор в формулах и сравнить суть функции на бесконечности.

\paragraph{Критерий сходимости}
Ряд $\sum_{n=1}^\infty a_n$ с неотрицательными членами ($a_n \ge 0$) сходится тогда и только тогда, когда последовательность его частичных сумм $S_n$ ограничена сверху: $\exists M > 0, \forall n : S_n \le M$.

\begin{proof}[Доказательство критерия]
Так как $a_n \ge 0$, каждая следующая частичная сумма $S_{n+1} = S_n + a_{n+1} \ge S_n$. Значит, последовательность $\{S_n\}$ монотонно не убывает. По известной теореме Вейерштрасса, монотонно неубывающая последовательность сходится тогда и только тогда, когда она ограничена сверху.
\end{proof}

\paragraph{Признаки сравнения (в форме неравенств)}
Пусть даны два неотрицательных ряды $\sum a_n$ и $\sum b_n$, причем начиная с некоторого номера выполняется $0 \le a_n \le b_n$.
1. Если $\sum b_n$ (мажорантный ряд) сходится, то сходится и $\sum a_n$.
2. Если $\sum a_n$ (минорантный ряд) расходится, то расходится и $\sum b_n$.

\begin{proof}
Пусть $S_n^A = \sum_{k=1}^n a_k$ и $S_n^B = \sum_{k=1}^n b_k$. Из неравенства следует $S_n^A \le S_n^B$.
1. Если $\sum b_n$ сходится к $B$, то $S_n^B \le B$. Тогда $S_n^A \le B$, то есть последовательность $S_n^A$ ограничена сверху. По критерию выше — ряд $\sum a_n$ сходится.
2. Если $\sum a_n$ расходится, то $S_n^A \to +\infty$. Так как $S_n^B \ge S_n^A$, по теореме о двух милиционерах (о зажатой последовательности) $S_n^B \to +\infty$, то есть ряд расходится.
\end{proof}

\paragraph{Предельный признак сравнения}
Пусть $a_n > 0$ и $b_n > 0$. Если существует конечный и отличный от нуля предел $\lim_{n \to \infty} \frac{a_n}{b_n} = c \in (0, +\infty)$, то ряды $\sum a_n$ и $\sum b_n$ сходятся или расходятся одновременно.

\begin{proof}
По определению предела, для $\varepsilon = \frac{c}{2}$ найдется номер $N$, что для $n > N$ выполняется:
$$ c - \frac{c}{2} < \frac{a_n}{b_n} < c + \frac{c}{2} \implies \frac{c}{2} b_n < a_n < \frac{3c}{2} b_n $$
Получили двустороннее неравенство. 
Левое неравенство $\left(\frac{c}{2} b_n < a_n\right)$ гарантирует сходимость $\sum b_n$ при сходимости $\sum a_n$ (или расходимость $a_n$ при расходимости $b_n$). 
Правое неравенство $\left(a_n < \frac{3c}{2} b_n\right)$ гарантирует сходимость $\sum a_n$ при сходимости $\sum b_n$. 
Ряды ведут себя абсолютно одинаково.
\end{proof}

\paragraph{Пример с тригонометрией}
Исследуем на сходимость ряд $\sum_{n=1}^\infty \sin\left(\frac{1}{n^2}\right)$. 
Знаем, что при $x \to 0$ выполняется эквивалентность $\sin x \sim x$. При $n \to \infty$ аргумент $\frac{1}{n^2} \to 0$. Сравним в предельной форме с эталонным обобщенным гармоническим рядом $\sum \frac{1}{n^2}$:
$$ \lim_{n \to \infty} \frac{\sin(1/n^2)}{1/n^2} = 1 \in (0, +\infty) $$
Так как эталонный ряд $\sum \frac{1}{n^2}$ сходится (степень $p=2 > 1$), то по предельному признаку сравнения сходится и исходный ряд с синусом.

Что ж, перейдем к настоящим легендам рядов, гениям и просто уважаемым людям, а точнее к их признакам.


% --- Для Вопроса 32 (Ряды с неотрицательными членами) ---
\begin{tcolorbox}[colback=gray!5, colframe=black, boxrule=0.5pt, arc=0pt]
\textbf{Резюмируем:} \\
\textbf{Тезис:} $a_n \ge 0$ сходится $\iff$ частичные суммы ограничены ($S_n \le M$). \\
\textbf{Ход (Неравенства):} Пусть $a_n \le b_n$, тогда $S_n^A \le S_n^B$. Если $B$ сходится, то $S_n^B \le B \implies S_n^A \le B$ (ограничена, сходится). \\
\textbf{Ход (Предельный):} $\lim (a_n/b_n) = c > 0$. По пределу $\frac{c}{2} b_n < a_n < \frac{3c}{2} b_n$. Ряды зажаты и ведут себя одинаково.
\end{tcolorbox}


% --- Вопрос 33 ---
\examquestion{33}{Признак Д'Аламбера}
{Сформулируйте признак Д'Аламбера. Проведите доказательство признака, используя метод мажорирования исследуемого ряда сходящейся или расходящейся геометрической прогрессией. Обоснуйте, почему при $q=1$ признак не работает.}

\paragraph{Суть на пальцах}
Д'Аламбер смотрит, во сколько раз следующий шаг меньше предыдущего. Если каждый раз число строго уменьшается на один и тот же гарантированный процент (образуя убывающую геометрическую прогрессию), то ряд сходится. Если коэффициент застрял на единице — тест слепнет, потому что скорость убывания слишком пограничная.

\paragraph{Формулировка}
Пусть для ряда с положительными членами $\sum_{n=1}^\infty a_n$ существует предел отношения последующего члена к предыдущему: $\lim_{n \to \infty} \frac{a_{n+1}}{a_n} = q$.
Тогда:
1. Если $q < 1$, ряд сходится.
2. Если $q > 1$, ряд расходится.

\begin{proof}[Доказательство (через геометрическую прогрессию)]
1. \textbf{Случай $q < 1$:} Выберем число $r$, лежащее строго между пределом и единицей: $q < r < 1$. По определению предела, найдется номер $N$, начиная с которого все отношения $\frac{a_{n+1}}{a_n} < r$. 
Тогда: $a_{N+1} < a_N r$, $a_{N+2} < a_{N+1} r < a_N r^2$, и в общем виде:
$$ a_{N+k} < a_N r^k $$
Мы ограничили хвост нашего ряда (начиная с $N$) членами геометрической прогрессии $\sum a_N r^k$. Так как знаменатель прогрессии $r < 1$, эта прогрессия сходится. По признаку сравнения в форме неравенств, наш ряд тоже сходится.

2. \textbf{Случай $q > 1$:} По определению предела, начиная с некоторого $N$, выполняется $\frac{a_{n+1}}{a_n} \ge 1$. Это значит, что $a_{n+1} \ge a_n > 0$. Последовательность членов ряда монотонно возрастает и не может стремиться к нулю. Нарушен необходимый признак сходимости ($\lim a_n \neq 0$), следовательно, ряд расходится.
\end{proof}

\paragraph{Почему признак не работает при $q=1$}
При $q=1$ предел не дает достаточно информации о скорости убывания членов. Рассмотрим два классических ряда:
* Гармонический ряд $\sum \frac{1}{n}$: $\lim \frac{1/(n+1)}{1/n} = \lim \frac{n}{n+1} = 1$. Ряд \textbf{расходится}.
* Обобщенный гармонический $\sum \frac{1}{n^2}$: $\lim \frac{1/(n+1)^2}{1/n^2} = \lim \left(\frac{n}{n+1}\right)^2 = 1$. Ряд \textbf{сходится}.
В обоих случаях предел Д'Аламбера равен 1, но судьба рядов противоположна. Для таких случаев нужны более тонкие признаки.

% --- Для Вопроса 33 (Признак Д'Аламбера) ---
\begin{tcolorbox}[colback=gray!5, colframe=black, boxrule=0.5pt, arc=0pt]
\textbf{Резюмируем:} \\
\textbf{Тезис:} $\lim (a_{n+1}/a_n) = q$. $q<1$ сходится, $q>1$ расходится. $q=1$ ломается. \\
\textbf{Ход ($q<1$):} Берем $q < r < 1$. С некоторого $N$ верно $a_{n+1}/a_n < r$. Хвост мажорируется сходящейся геометрической прогрессией $a_N r^k$. \\
\textbf{Ход ($q>1$):} $a_{n+1} \ge a_n > 0$. Общий член не стремится к нулю, ряд расходится. При $q=1$ тесты на $\sum 1/n$ и $\sum 1/n^2$ дают разные итоги.
\end{tcolorbox}

% --- Вопрос 34 ---
\examquestion{34}{Радикальный признак Коши}
{Сформулируйте радикальный признак Коши. Проведите его доказательство. Проведите сравнительный анализ «силы» признаков Д'Аламбера и Коши: докажите, что из существования предела отношения членов $\lim(a_{n+1}/a_{n})=q$ всегда следует существование предела корней $\lim\sqrt[n]{a_{n}}=q$, но не наоборот.}

\paragraph{Суть на пальцах}
Признак Коши извлекает корень $n$-й степени из $n$-го элемента. Идейно он делает то же самое, что и Д'Аламбер — ищет скрытую геометрическую прогрессию. Но математически он более «всеядный». Если члены ряда скачут то вверх, то вниз, Д'Аламбер сойдет с ума (предел отношения не будет существовать), а Коши жестким корнем усреднит эти колебания и выдаст точный вердикт.

\paragraph{Формулировка и доказательство признака Коши}
Пусть для ряда с неотрицательными членами существует предел $\lim_{n \to \infty} \sqrt[n]{a_n} = q$.
Если $q < 1$ — ряд сходится. Если $q > 1$ — ряд расходится. (При $q=1$ признак не дает ответа).

\begin{proof}
1. \textbf{Случай $q < 1$:} Выберем $q < r < 1$. Начиная с некоторого номера $N$, выполняется $\sqrt[n]{a_n} < r$. Возводя в $n$-ю степень, получаем $a_n < r^n$. Хвост нашего ряда ограничен сверху сходящейся геометрической прогрессией ($r < 1$). Значит, исходный ряд сходится.
2. \textbf{Случай $q > 1$:} Начиная с некоторого $N$, выполняется $\sqrt[n]{a_n} \ge 1 \implies a_n \ge 1$. Общий член не стремится к нулю, необходимый признак нарушен, ряд расходится.
\end{proof}

\paragraph{Анализ «силы»: Коши мощнее Д'Аламбера}
\textbf{Теорема связи:} Если последовательность положительных чисел $a_n$ такова, что существует $\lim_{n \to \infty} \frac{a_{n+1}}{a_n} = q$, то существует и $\lim_{n \to \infty} \sqrt[n]{a_n} = q$.

\begin{proof}[Доказательство теоремы связи]
Пусть $\lim_{n \to \infty} \frac{a_{n+1}}{a_n} = q$. По определению предела, для любого $\varepsilon > 0$ найдется номер $N$, что для всех $k \ge N$ выполняется:
$$ q - \varepsilon < \frac{a_{k+1}}{a_k} < q + \varepsilon $$
Распишем член $a_n$ (для $n > N$) в виде телескопического произведения дробей:
$$ a_n = a_N \cdot \frac{a_{N+1}}{a_N} \cdot \frac{a_{N+2}}{a_{N+1}} \cdots \frac{a_n}{a_{n-1}} $$
Здесь ровно $(n-N)$ дробей, и каждая из них зажата между $(q - \varepsilon)$ и $(q + \varepsilon)$. Получаем строгую оценку для $a_n$:
$$ a_N (q - \varepsilon)^{n-N} < a_n < a_N (q + \varepsilon)^{n-N} $$
Извлечем корень $n$-й степени из всех частей неравенства:
$$ \sqrt[n]{a_N (q - \varepsilon)^{-N}} \cdot (q - \varepsilon) < \sqrt[n]{a_n} < \sqrt[n]{a_N (q + \varepsilon)^{-N}} \cdot (q + \varepsilon) $$
Обозначим константы $C_1 = a_N (q - \varepsilon)^{-N}$ и $C_2 = a_N (q + \varepsilon)^{-N}$. Известно, что корень $n$-й степени из любой положительной константы стремится к единице при $n \to \infty$. (Чтобы Лариса Исхаквна точно одобрила, можно логарифмировать и по Лопиталю). Следовательно, $\sqrt[n]{C_1} \to 1$ и $\sqrt[n]{C_2} \to 1$. 
Переходя к пределу, получаем, что $\sqrt[n]{a_n}$ зажата между $(q - \varepsilon)$ и $(q + \varepsilon)$. В силу произвольности $\varepsilon$, по теореме о зажатой переменной $\lim_{n \to \infty} \sqrt[n]{a_n} = q$.
\end{proof}

\textbf{Контрпример (почему наоборот не работает):}
Построим ряд, где Коши работает, а Д'Аламбер ломается. Пусть ряд задан так: $a_n = 2^{-n}$ при четном $n$, и $a_n = 2^{-n+2}$ при нечетном $n$.
Посмотрим на предел Д'Аламбера $\frac{a_{n+1}}{a_n}$:
* Переход от нечетного к четному: $\frac{2^{-(n+1)}}{2^{-n+2}} = 2^{-3} = \frac{1}{8}$
* Переход от четного к нечетному: $\frac{2^{-(n+1)+2}}{2^{-n}} = 2^1 = 2$
Предел отношения скачет между $1/8$ и $2$ — Д'Аламбер \textbf{не существует}.
А теперь применим корень Коши: $\lim \sqrt[n]{a_n}$. Корень $n$-й степени из $2^{-n}$ равен $1/2$, а корень из $2^{-n+2} = 2^{2/n}/2 \to 1/2$. В обоих случаях предел равен $\frac{1}{2} < 1$. Коши уверенно заявляет: ряд \textbf{сходится}.



% --- Для Вопроса 34 (Радикальный признак Коши) ---
\begin{tcolorbox}[colback=gray!5, colframe=black, boxrule=0.5pt, arc=0pt]
\textbf{Резюмируем:} \\
\textbf{Тезис:} $\lim \sqrt[n]{a_n} = q$. $q<1$ сходится, $q>1$ расходится. Коши мощнее Д'Аламбера. \\
\textbf{Ход ($q<1$):} Выбираем $q < r < 1$. С некоторого $N$ верно $\sqrt[n]{a_n} < r \implies a_n < r^n$. Сравнили со сходящейся прогрессией. \\
\textbf{Сила:} Если есть предел Д'Аламбера, то Коши выдаст то же самое. Обратно неверно: для чередующегося $2^{-n}$ и $2^{-n+2}$ Д'Аламбер скачет, а Коши четко дает $1/2$.
\end{tcolorbox}


% --- Вопрос 35 ---
\examquestion{35}{Абсолютная сходимость и ее свойства}
{Сформулируйте и докажите теорему о том, что из абсолютной сходимости числового ряда следует его обычная сходимость. Примените критерий Коши и неравенство треугольника. Сформулируйте теорему Римана о перестановке членов условно сходящегося ряда, приведя примеры.}

\paragraph{Определения}
Ряд $\sum_{n=1}^\infty a_n$ называется \textbf{абсолютно сходящимся}, если сходится ряд из абсолютных величин его членов $\sum_{n=1}^\infty |a_n|$.
Ряд называется \textbf{условно сходящимся}, если сам ряд $\sum_{n=1}^\infty a_n$ сходится, а ряд из абсолютных величин $\sum_{n=1}^\infty |a_n|$ расходится.

\paragraph{Теорема}
Если ряд $\sum_{n=1}^\infty |a_n|$ сходится, то ряд $\sum_{n=1}^\infty a_n$ также сходится.

\begin{proof}[Доказательство]
1. Так как ряд $\sum |a_n|$ сходится, по критерию Коши для любого $\varepsilon > 0$ существует номер $N$ такой, что для любого $n > N$ и любого натурального $p$ выполняется:
$$ \sum_{k=n+1}^{n+p} |a_k| < \varepsilon $$
2. Рассмотрим частичную сумму «хвоста» исходного ряда $\sum a_n$ и применим неравенство треугольника:
$$ \left| \sum_{k=n+1}^{n+p} a_k \right| \le \sum_{k=n+1}^{n+p} |a_k| $$
3. Учитывая условие из пункта 1, получаем:
$$ \left| \sum_{k=n+1}^{n+p} a_k \right| < \varepsilon $$
4. Согласно критерию Коши, выполнение этого условия означает, что ряд $\sum_{n=1}^\infty a_n$ сходится.
\end{proof}

\paragraph{Признак Лейбница (для условной сходимости)}
Чтобы ряд сходился условно, он часто должен быть знакочередующимся. Признак Лейбница гласит:
Если для знакочередующегося ряда $\sum (-1)^{n+1} b_n$ (где $b_n > 0$) выполняются два условия:
1. Последовательность $\{b_n\}$ монотонно убывает: $b_{n+1} \le b_n$.
2. Предел общего члена равен нулю: $\lim_{n \to \infty} b_n = 0$.
Тогда ряд сходится.

\textbf{Пример:} Ряд $\sum (-1)^{n+1} \frac{1}{n} = 1 - \frac{1}{2} + \frac{1}{3} - \frac{1}{4} + \dots$
Здесь $b_n = 1/n$. Условия Лейбница выполнены ($1/n$ убывает и стремится к 0), значит, ряд сходится. Но ряд из модулей $\sum 1/n$ — это гармонический ряд, он расходится. Значит, наш ряд сходится \textbf{условно}.

\paragraph{Теорема Римана о перестановке членов}
Если ряд сходится условно, то путем перестановки его членов можно получить ряд, сходящийся к любому наперед заданному числу $S$, или ряд, расходящийся к $\pm\infty$.

\textbf{Пример:} Возьмем тот же знакочередующийся гармонический ряд:
$$ 1 - \frac{1}{2} + \frac{1}{3} - \frac{1}{4} + \frac{1}{5} - \frac{1}{6} + \dots = \ln 2 $$
Если мы переставим члены так, чтобы брать один положительный и два отрицательных (например: $1 - \frac{1}{2} - \frac{1}{4} + \frac{1}{3} - \frac{1}{6} - \frac{1}{8} + \dots$), то сумма этого ряда изменится и станет равна $\frac{1}{2} \ln 2$. 

Этот парадокс возможен только для условно сходящихся рядов, потому что «положительная» и «отрицательная» части таких рядов по отдельности бесконечно велики, и меняя их порядок, мы меняем баланс бесконечностей. В абсолютно сходящихся рядах такой «магии» нет — сумма всегда инвариантна относительно перестановки.


% --- Для Вопроса 35 (Абсолютная сходимость) ---
\begin{tcolorbox}[colback=gray!5, colframe=black, boxrule=0.5pt, arc=0pt]
\textbf{Резюмируем:} \\
\textbf{Тезис:} Из сходимости $\sum |a_n|$ следует сходимость $\sum a_n$. \\
\textbf{Ход:} По критерию Коши для ряда модулей: $\sum_{k=n+1}^{n+p} |a_k| < \varepsilon$. \\
\textbf{Трюк:} По неравенству треугольника модуль суммы меньше суммы модулей: $|\sum a_k| \le \sum |a_k| < \varepsilon$. Исходный ряд прошел критерий Коши. \\
\textbf{Риман:} Члены абсолютно сходящегося ряда можно менять местами без потери суммы. Условно сходящийся ряд можно перетасовать под любую сумму.
\end{tcolorbox}